% ============================================================
% SKELETON -- Part II, Section 9: Thermodynamic Processes
% ============================================================

\section{Thermodynamic Processes}

\paragraph{Scope of this section.}
Sections~6--8 describe equilibrium states. This section connects the states
by processes, computes the heat and the work along each process, and
assembles the processes into cycles. The ideal gas of Section~7 and the van
der Waals gas of Section~8 serve as the worked examples.

\subsection{Quasi-Static Processes}

\paragraph{Processes through equilibrium states.}
A \textcolor{blue}{quasi-static process} passes through equilibrium states
only, so the state variables are defined along the whole path. Heat and
work then follow from integrals along the path, and both depend on the path.
% FROM: delta Q, delta W and the Clausius inequality (Section 3).
\begin{equation}
    \boxed{W=\int P\,dV,\qquad Q=\int T\,dS\quad\text{(reversible)}.}
\end{equation}

\paragraph{Reversible and irreversible processes.}
A \textcolor{blue}{reversible process} is quasi-static and produces no
entropy, so the reverse path restores the system and the surroundings. A
process with friction, heat flow across a finite temperature difference or
free expansion produces entropy and is irreversible.

\subsection{Elementary Processes}

\paragraph{Isothermal, isobaric, isochoric and adiabatic processes.}
Each elementary process fixes one variable. The heat and the work of each
process follow from the response functions of Section~6.
Table~\ref{tab:processes} summarizes the four processes for an ideal gas.
% TODO: general expressions for delta Q and delta W; ideal-gas values.
\begin{table}[h]
\centering
\caption{Work done by an ideal gas and heat absorbed in the four elementary
processes between states 1 and 2.}
\label{tab:processes}
\begin{tabular}{lccc}
    Process & Fixed variable & $W$ & $Q$ \\ \hline
    Isothermal & $T$ & $Nk_BT\ln(V_2/V_1)$ & $W$ \\
    Isobaric & $P$ & $P(V_2-V_1)$ & $C_P(T_2-T_1)$ \\
    Isochoric & $V$ & $0$ & $C_V(T_2-T_1)$ \\
    Adiabatic & $S$ & $C_V(T_1-T_2)$ & $0$ \\
\end{tabular}
\end{table}

\paragraph{Adiabats and isotherms.}
An adiabat is steeper than an isotherm in the $P$--$V$ plane by the factor
$\gamma\equiv C_P/C_V$. For an ideal gas with constant $\gamma$, the
adiabat is a power law.
% FROM: C_P / C_V = kappa_T / kappa_S (Section 6).
\begin{equation}
    \boxed{\left(\frac{\partial P}{\partial V}\right)_{\!S}
    =\gamma\left(\frac{\partial P}{\partial V}\right)_{\!T},
    \qquad PV^\gamma=\text{const}\quad\text{(ideal gas)}.}
\end{equation}

\subsection{Heat Engines and the Carnot Cycle}

\paragraph{Cycles, engines and efficiency.}
A \textcolor{blue}{cycle} returns the system to the initial state, so
$\Delta U=0$ and the net work equals the net heat. A
\textcolor{blue}{heat engine} absorbs the heat $Q_{\text{h}}$ from a hot
reservoir, rejects the heat $Q_{\text{l}}$ to a cold reservoir and
delivers the work $W=Q_{\text{h}}-Q_{\text{l}}$.
\textcolor{red}{Here, the subscripts h and l label the hot and the cold
reservoir, so that the subscript c stays reserved for the critical point.}
\begin{equation}
    \boxed{\eta\equiv\frac{W}{Q_{\text{h}}}=1-\frac{Q_{\text{l}}}{Q_{\text{h}}}.}
\end{equation}

\paragraph{The Carnot cycle.}
The \textcolor{blue}{Carnot cycle} consists of two isotherms at
$T_{\text{h}}$ and $T_{\text{l}}$ joined by two adiabats. The cycle
exchanges heat with each reservoir at a single temperature, so the cycle is
reversible.
% TODO: four steps in the P--V and T--S planes (figures); heat and work of
%       each step for an ideal gas; V_2/V_1 = V_3/V_4 from the adiabats.
\begin{equation}
    \boxed{\frac{Q_{\text{l}}}{Q_{\text{h}}}=\frac{T_{\text{l}}}{T_{\text{h}}},
    \qquad
    \eta_{\text{C}}=1-\frac{T_{\text{l}}}{T_{\text{h}}}.}
\end{equation}

\paragraph{The Carnot bound from the second law.}
The Carnot efficiency does not depend on the working substance. Over any
cycle between two reservoirs, the Clausius inequality gives
$Q_{\text{h}}/T_{\text{h}}-Q_{\text{l}}/T_{\text{l}}\leq0$, so no engine
exceeds the Carnot efficiency, and every reversible engine reaches the
Carnot efficiency.
% FROM: Clausius inequality (Section 3).
\begin{equation}
    \boxed{\eta\leq\eta_{\text{C}},\qquad
    \oint\frac{\delta Q}{T}\leq0.}
\end{equation}

\paragraph{Refrigerators and heat pumps.}
A Carnot cycle run backward moves heat from the cold to the hot reservoir
with the work $W$. The \textcolor{blue}{coefficient of performance} measures
the useful heat per unit work.
\begin{equation}
    \boxed{\mathrm{COP}_{\text{ref}}=\frac{Q_{\text{l}}}{W}\leq\frac{T_{\text{l}}}{T_{\text{h}}-T_{\text{l}}},
    \qquad
    \mathrm{COP}_{\text{hp}}=\frac{Q_{\text{h}}}{W}\leq\frac{T_{\text{h}}}{T_{\text{h}}-T_{\text{l}}}.}
\end{equation}

\paragraph{Thermodynamic temperature.}
The ratio $Q_{\text{l}}/Q_{\text{h}}$ of a reversible engine depends on the
two reservoir temperatures only. This ratio defines a temperature scale
independent of any substance, and the scale coincides with the $T$ of
Section~3.
% TODO: link to beta = 1/(k_B T) as an empirical correspondence (Section 3).

\paragraph{Other cycles.}
Real engines use other cycles. Each cycle combines the elementary processes
and has an efficiency below the Carnot bound.
% TODO: Otto (two adiabats, two isochores) and Stirling cycles as exercises.

\subsection{Free Expansion and Throttling}

\paragraph{Free expansion.}
In a \textcolor{blue}{free expansion}, a gas expands into vacuum inside an
insulated container, so neither heat nor work is exchanged and $U$ stays
constant. The process is irreversible, and the temperature change measures
the deviation from ideal behavior.
% TODO: derive with the second Maxwell relation; Delta S = N k_B ln(V_2/V_1)
%       for an ideal gas; cooling of a van der Waals gas.
\begin{equation}
    \boxed{\left(\frac{\partial T}{\partial V}\right)_{\!U}
    =-\frac{1}{C_V}\left[T\left(\frac{\partial P}{\partial T}\right)_{\!V}-P\right].}
\end{equation}

\paragraph{Throttling.}
In a \textcolor{blue}{throttling process}, a gas flows steadily through a
porous plug from high to low pressure. The enthalpy stays constant, and the
\textcolor{blue}{Joule--Thomson coefficient} gives the temperature change.
The coefficient vanishes for an ideal gas.
% TODO: H = const from the work balance on the two sides of the plug;
%       derive with the fourth Maxwell relation.
\begin{equation}
    \boxed{\mu_{\text{JT}}\equiv\left(\frac{\partial T}{\partial P}\right)_{\!H}
    =\frac{V}{C_P}\left(T\alpha_P-1\right).}
\end{equation}

\paragraph{Joule--Thomson inversion.}
The attraction cools a throttled gas, and the hard core heats the gas. At
low density, the van der Waals gas of Section~8 cools on throttling below
the \textcolor{blue}{inversion temperature}. Gas liquefaction relies on
this cooling.
% TODO: mu_JT ~ (2a/(k_B T) - b)/c_P from B_2.
\begin{equation}
    \boxed{k_BT_{\text{inv}}=\frac{2a}{b}=\frac{27}{4}\,k_BT_c.}
\end{equation}
